\manualchapter{Control reference}{sec:controls}
Panel map: \textbf{1} global controls;
\textbf{2--5} operators 1--4.

The left block selects algorithm, feedback, LFO and output saturation. Four
repeated operator blocks contain envelope knobs beside tuning and modulation
controls; each column has its own CV, gate, retrigger and pitch sockets.

Algorithm spans 0--7 (default 7); feedback spans 0--7 (default 0) and feeds
operator 1. LFO spans 0--7 (default 0). Their CV inputs add $7V/8$ to the
knob, truncate, and clamp. A 1\,V algorithm step is therefore not a reliable
one-index step. The current diagram display uses a different CV scale from
the audio engine; use the knob without CV when learning the routings.

SAT spans 0--127 (default 127). Its CV adds $127V/10$, clamped to 0--127.
Lower values attenuate the operator sum before 14-bit clipping; SAT does not
invert polarity. Both outputs carry the same approximately $\pm5$\,V signal.

\subsection{Envelope knobs in each operator block}
\begin{tabularx}{\textwidth}{@{}l l X@{}}
\toprule Control & Range; default & Result \\\midrule
AR & 1--31; 31 & Attack rate; larger values rise faster. \\
TL & 0--100; 100 & Peak level; larger values are louder. \\
D1R & 0--31; 0 & First decay rate; zero holds the peak. \\
SL / T1L & 0--15; 15 & Level at which the second decay begins. \\
D2R & 0--31; 0 & Second decay rate; zero holds this level. \\
RR & 0--15; 15 & Release rate after gate falls. \\
\bottomrule
\end{tabularx}

These are chip rate and attenuation settings, not seconds or linear amplitude
percentages. Rate scaling (0--3, default 0) makes envelopes respond more quickly
at higher pitches. Loop is off by default; it enables the chip's SSG envelope
mode. Use it for repeating contours after first establishing a gated sound.
The hardware envelope terminology is retained in the reference figure and
in \citet{genesis-sound-software-manual}.

Each envelope CV adds $mV/8$ to its knob, where $m$ is that knob's maximum.
The sum is truncated to an integer and clamped to the control range. For
example, with AR at 15, $+4$\,V adds 15.5 before truncation, giving 30.
This is an offset; a 0\,V cable leaves the slider's setting intact.
